\documentclass[10pt,a4paper]{amsart}
\usepackage[margin=19mm]{geometry}
\usepackage{amsmath,amssymb,mathtools}
\usepackage[scr=rsfs,cal=euler]{mathalfa}
\usepackage{xfrac}
\usepackage[unicode,hidelinks,bookmarks=false]{hyperref}
\newcommand{\ie}{i.e.\ }
\newcommand{\cf}{cf.\ }
\newcommand{\resp}{respectively\ }
\newcommand{\Subsection}{\subsection}
\providecommand{\nobreakdash}{\nobreak\hbox{-}}
\setlength{\emergencystretch}{2em}
\multlinegap0pt
\usepackage{braket}
\usepackage{amssymb}
\usepackage{dsfont}
\usepackage{xcolor}
\newtheorem{theorem}{Theorem}
\newcommand{\cS}{\mathcal{S}}
\title{Rate-Reliability Tradeoff for Deterministic Identification over Gaussian Channels}
\author{Pau Colomer, Christian Deppe, Holger Boche, Andreas Winter}
\date{Source: arXiv:2602.12182v1 (CC BY 4.0)}
\begin{document}
\maketitle

\noindent\emph{Source and licence.} Rate-Reliability Tradeoff for Deterministic Identification over Gaussian Channels. Version arXiv:2602.12182v1 (CC BY 4.0), distributed by arXiv under the Creative Commons Attribution 4.0 International licence, http://creativecommons.org/licenses/by/4.0/, as recorded in the arXiv licence metadata for this exact version and shown on the abstract page https://arxiv.org/abs/2602.12182v1. Full licence notice: \texttt{licenses/arxiv-2602.12182-CC-BY-4.0.txt}.\par\smallskip

\noindent\textit{Editorial scope.} Counted unit: the article's theorem thm:achievability\_linearithmic (source0.tex lines 694--700). The article's complete original proof of it is delivered with it (source0.tex lines 702--715, one \texttt{proof} environment of 14 lines). No mathematics is written, repaired or re-derived: the statement and the proof are the article's own lines, in the article's own notation, and no retained line is substituted or re-flowed.  Retained uncounted as the source-local context the proof needs: lines 561--568; lines 659--661; lines 668--674.  Nothing was omitted from any retained span, so the package carries no omission record.  No retained line contains a citation, so the wrapper carries no bibliography and no bibliography entry.  No retained line contains a figure environment, an image-inclusion command or drawing code, so no image is delivered and the package declares no asset dependency.  Editorial formatting only, in addition: an AMS \texttt{amsart} preamble that restates the article's own declarations and macros used by the retained lines, namely the environments \texttt{theorem} on separate counters, together with the standard packages those lines need.  The retained environments print in this document as Theorem 1 in order of appearance, whereas the article prints the same items with the numbering of its own full document.  The complete native source of the article as distributed in the arXiv e-print for this exact version is bundled unchanged as \texttt{sources/194527/source0.tex}.
\begin{equation}
\label{eq:code_size}
\begin{split}
N&=\frac{\textsc{Vol}\left[\bigcup_{i=1}^N\cS_i(n,r)\right]}{\textsc{Vol}\left[\cS_i(n,r)\right]}\\
&\geq\frac{\Delta_n \textsc{Vol}\left[\cS_-(n,\sqrt{nP}-r)\right]}{\textsc{Vol}\left[\cS_i(n,r)\right]}\\
&\geq2^{-n}\left[\frac{\sqrt{nP}-r}{r}\right]^n=\left(\frac{1-\epsilon}{2\epsilon}\right)^n,
\end{split}
\end{equation}

\begin{equation}\label{eq:condition}
P\epsilon(n)^2>\nu_M \sqrt{E_1(n)}.
\end{equation}

\begin{theorem}\label{thm:ach_linear}
Let $\nu_M$ be the maximum eigenvalue of the covariance matrix $\Sigma$, and $P$ the maximum power constraint of the linear Gaussian channel. Then, for any arbitrarily small $\tau>0$ we can construct a DI code with linear rate lower bounded as:
\[
R>\log\left(\sqrt{\frac{P}{4\nu_M(1+\tau) \sqrt{E_1}}}-\frac12\right),
\]
and exponentially fast vanishing errors $\lambda_1:=e^{-nE_1}$ with $0<\sqrt{E_1}< \min\{1,\frac{P}{9\tau\nu_M}\}$, and $\lambda_2:=e^{-nE_2}$ with $E_2=\frac{\tau^2E_1}{1+P/\nu_M}$.
\end{theorem}

\begin{theorem}\label{thm:achievability_linearithmic}
Given $\tau>0$ and $0<\beta<1$ we can construct a DI code of linearithmic size with linearithmic rate
\[
\frac{R(n)}{\log n}:=\frac{\log N}{n\log n}\geq\frac{\beta}{4}+O\left(\frac{1}{\log n}\right),
\]
and sub-exponentially fast vanishing errors $\lambda_1:=e^{-nE_1(n)}$ and $\lambda_2:=e^{-nE_2(n)}$, with $E_1(n)=n^{-\beta}$ and $E_2(n)=\frac{2\tau^2n}{1+P/\nu_M}n^{-\beta}$.
\end{theorem}

\begin{proof}
Let us start by defining a polynomially-fast vanishing first type error exponent $E_1(n):=n^{-\beta}$ for $0<\beta<1$, which implies a sub-exponentially fast vanishing first type of error, and choosing $\epsilon(n)^2=(1+\tau)\nu_M\sqrt{E_1(n)}/P$ such that the condition in Eq.~\eqref{eq:condition} is respected. Notice that it is the same $\epsilon$ we chose in the proof of Theorem \ref{thm:ach_linear} but with the new definition of $E_1(n)=n^{-\beta}$, therefore, the exponent of the second type of error is also equal $E_2=\frac{2\tau^2E_1(n)}{1+P/\nu_M}$. It is only left to calculate the size of the code through Eq.~\eqref{eq:code_size}:
\[
N\geq\left(\frac{n^{\frac{\beta}{4}}\sqrt{P}}{2\sqrt{(1+\tau)\nu_M}}-\frac12\right)^n\geq\left(\frac{n^{\frac{\beta}{4}}\sqrt{P}}{4\sqrt{(1+\tau)\nu_M}}\right)^n.
\]
Then, the linear rate can be written as 
\[
\begin{split}
R:=\frac{\log N}{n}&\geq\log\left(\frac{n^{\frac{\beta}{4}}\sqrt{P}}{4\sqrt{(1+\tau)\nu_M}}\right)\\   
&=\frac{\beta}{4}\log(n)+\frac12\log\left(\frac{P}{16(1+\tau)\nu_M}\right).
\end{split}
\]
In other words, the linearithmic rate is $\frac{\log N}{n\log n}\geq\frac{\beta}{4}+O(\frac{1}{\log n})$. The proof is completed.
\end{proof}

\end{document}
